Let $P^{n}(c)$ be the set of collections $(A_1,A_2;C_1,\dots,C_{n};e)$
in $\operatorname{End}(\mathbb{C}^{c})^{\oplus n+2}\oplus\operatorname{Hom}(\mathbb{C}^{c},\mathbb{C})$ sa\-tisfying the conditions\\[8pt]
(P1)\ \ \
$\displaystyle
\begin{cases}
A_1C_1A_{2}=A_2C_{1}A_{1},&\qquad\text{when $n=1$},\\[2pt]
\begin{aligned}
A_1C_q&=A_2C_{q+1},\\
C_qA_1&=C_{q+1}A_2
\end{aligned}
\qquad\text{for}\quad q=1,\dots,n-1,&\qquad\text{when $n>1$;}
\end{cases}
$
\\[8pt]
(P2) \
$A_1+\lambda A_2$ is a \emph{regular pencil} of matrices, i.e., there exists $[\nu_{1},\nu_{2}]\in\mathbb{P}^1$ such that $\det(\nu_1A_1+\nu_2A_2)\neq0$;
\\[8pt]
(P3) \
for all values of the parameters $ ([\lambda_1,\lambda_2],(\mu_1,\mu_{2}) )\in\mathbb{P}^1\times\mathbb{C}^{2}$ satisfying
\begin{equation*}
\lambda_{1}^{n}\mu_{1}+\lambda_{2}^{n}\mu_{2}=0
\end{equation*}
there is no nonzero vector $v\in\mathbb{C}^c$ such that
\begin{equation*}
\begin{cases}
C_{1}A_{2}v=-\mu_1v,\\
C_{n}A_{1}v=(-1)^n\mu_2v,\\
v\in\ker e
\end{cases}
\qquad\text{and}\qquad\left(\lambda_2{A_1}+\lambda_1{A_2}\right)v=0 .
\end{equation*}

The group
$\operatorname{GL}(c, \mathbb{C})\times \operatorname{GL}(c, \mathbb{C})$ acts on $P^n(c)$ according to
\begin{equation*} (A_i, C_j, e ) \mapsto \big(\phi_2A_i\phi_1^{-1},\phi_1C_j\phi_2^{-1}, e\phi_1^{-1}\big)
%\label{eqrho}
\end{equation*}
for $i=1,2$, $j=1,\dots,n$, $ (\phi_1,\phi_2)\in\operatorname{GL}(c, \mathbb{C})\times \operatorname{GL}(c, \mathbb{C})$.

 The following result expresses the fact that the collections $(A_1,A_2;C_1,\dots,C_{n};e)$ satisfying
 conditions (P1) to (P3) are ADHM data for the varieties $\operatorname{Hilb}^c (X_{n})$ (this is Theorem~3.1 in~\cite{bbr}).
 \begin{Theorem}\label{thm0}
$P^{n}(c)$ is a principal $\operatorname{GL}(c, \mathbb{C})\times \operatorname{GL}(c, \mathbb{C})$-bundle over $\operatorname{Hilb}^c (X_{n})$.
\end{Theorem}

\section{The main result}\label{sectionmainresult}

\begin{figure}[t]\centering
$ \xymatrix@R-2.3em{
&&\mbox{\scriptsize$0$}&&\mbox{\scriptsize$1$}\\
&&\bullet\ar@/_3ex/[llddddddd]_{j}\ar@/^/[rr]^{a_{1}}\ar@/^4ex/[rr]^{a_{2}}&&\bullet\ar@/^/[ll]^{c_{1}}
\ar@/^4ex/[ll]^{c_2}\ar@{..}@/^8ex/[ll]\ar@{..}@/^9ex/[ll]
\ar@/^10ex/[ll]^{c_{n}}& \\ \\ \\ \\ \\&&&&& \\ \\
\bullet\ar@/_/[rruuuuuuu]^{i_1}\ar@/_3ex/[rruuuuuuu]_{i_2}\ar@{..}@/_6ex/[rruuuuuuu]\ar@{..}@/_8ex/[rruuuuuuu]\ar@/_10ex/[rruuuuuuu]_{i_{n-1}}&&&&&\\
\mbox{\scriptsize$\infty$}&&&&
}
$
\caption{The quivers $Q_n$.}\label{Qn}
\end{figure}
Now we turn to the purpose of this paper, namely, proving that the Hilbert schemes of points of the varieties $X_n$ are
isomorphic to moduli spaces of representations of suitable quivers.
 For any $n\geq2$ let $Q_{n}$ be the framed quiver in Fig.~\ref{Qn}, where $\infty$ is the framing vertex. Let $J_{n}$ be the two sided ideal of $\mathbb{C} Q_{n}$ generated by the relations
\begin{equation}
\begin{cases}
a_2c_{q+1}-a_1c_q =0,\\
c_{q+1}a_2-c_qa_1-i_qj =0
\end{cases}
\qquad\text{for}\quad q=1,\dots,n-1 .
\label{eq-gen-J}
\end{equation}
Our purpose is to describe the spaces of representations of the quiver $Q_n$ with relations $J_{n}$, i.e., the spaces of representations of the quotient algebra $\Lambda_{n}=\mathbb{C} Q_{n}/J_{n}$.

We recall some basic definitions. Given $\vec{v} = (v_0, v_1)\in \mathbb{N}^2$ and $w\in \mathbb{N}$,
a {\it $(\vec{v},w)$-dimensional representation of $\Lambda_n$}
is the datum of a triple of $\mathbb{C}$-vector spaces $V_0$, $V_1$, $W$, with $\dim V_i = v_i$, $\dim W=w$, and of an element
$(A_1,A_2;C_1, \dots,C_n;e;f_1,\dots,f_{n-1})$ in
\[
\operatorname{Hom}_\mathbb{C}(V_0,V_1)^{\oplus2}\oplus\operatorname{Hom}_\mathbb{C}(V_1,V_0)^{\oplus n}\oplus\operatorname{Hom}_\mathbb{C}(V_0,W)\oplus\operatorname{Hom}_\mathbb{C}(W,V_0)^{\oplus n-1}
\]
 satisfying the relations determined by equations~\eqref{eq-gen-J}, namely
 \begin{equation}
\begin{cases}
A_2 C_{q+1}- A_1 C_q =0,\\
C_{q+1} A_2 - C_q A_1- f_q e =0
\end{cases}
\qquad\text{for}\quad q=1,\dots,n-1 .
\label{eq-gen-Q1}\tag{Q1}
\end{equation}

The space $\operatorname{Rep}(\Lambda_n ,\vec{v},w)$ of all $(\vec{v},w)$-dimensional representations of $\Lambda_n$
is an affine variety, on which the group $G_{\vec{v}}=\operatorname{GL}(v_0,\mathbb{C})\times\operatorname{GL}(v_1,\mathbb{C})$ acts by basis change. Indeed, we ignore the action of $\operatorname{GL}(w,\mathbb{C})$ on the vector space~$W$ attached to the framing vertex. As usual, to get a~well behaved quotient space one has to perform a GIT construction by introducing a suitable notion of stability. This was done by A.~King \cite{KiQ} and, in a slightly different way, by A.~Rudakov~\cite{Rud97}. In the case of a quiver with a framing vertex, the following definition can be shown to be equivalent to the King--Rudakov one \cite{blr, CrBo}.
\begin{Definition}\label{def:stabil}
 Fix $\vartheta\in\mathbb R^2$. A $(\vec{v},w)$-dimensional representation {$(V_0,V_1,W)$} of $\Lambda_n$ is said to be \emph{$\vartheta$-semistable} if, for any subrepresentation $S=(S_0,S_1)$ {$\subseteq (V_0,V_1)$}, one has:
\begin{gather}
 \text{if $S_0\subseteq \ker e$, then $\vartheta\cdot(\dim S_0,\dim S_1)\leq0$} ;\label{eq:stabnak1}\\[5pt]
 \text{if $S_0\supseteq \operatorname{Im} f_i$ for $i=1,\dots,n-1$, then $\vartheta\cdot(\dim S_0,\dim S_1)\leq\vartheta\cdot(v_0,v_1)$} .\label{eq:stabnak2}
\end{gather}
A $\vartheta$-semistable representation is \emph{$\vartheta$-stable} if a strict inequality holds in \eqref{eq:stabnak1} whenever $S\neq0$ and
in \eqref{eq:stabnak2} whenever $S\neq(V_0,V_1)$.
\end{Definition}
Let $\operatorname{Rep}(\Lambda_n, \vec{v}, w)^{\rm ss}_\vartheta$ be the subset of $\operatorname{Rep}(\Lambda_n,\vec{v},w)$ consisting of $\vartheta$-semistable representations. By \cite[Proposition 5.2]{KiQ}, the coarse moduli space of
$(\vec{v},w)$-dimensional $\vartheta$-semistable representations of $\Lambda_n$ is the GIT quotient
\[\mathcal{M} (\Lambda_n,\vec{v},w, \vartheta)=\operatorname{Rep}(\Lambda_n, \vec{v}, w)^{\rm ss}_\vartheta/\!/ G_{\vec{v}} .\]
It can be proved that the open subset $\mathcal{M}^{\rm s} (\Lambda_n,\vec{v},w, \vartheta) \subset \mathcal{M} (\Lambda_n,\vec{v},w, \vartheta)$
consisting of stable representations makes up a fine moduli space. Notice that, for quivers without a framing, this
holds only when the dimension vector is primitive \cite[Proposition~5.3]{KiQ}, whilst this requirement is not necessary in the case of framed quivers~\cite{CrBo}. Theorem~4.5 of~\cite{bblr} states that the Hilbert scheme of points
$\operatorname{Hilb}^c (X_n)$ can be embedded into $\mathcal{M} (\Lambda_n,\vec{v},w, \vartheta)$
for suitable choices of $\vec{v}$, $w$, and $\vartheta$. Precisely, one has the following result:
\begin{Theorem}\label{thm45}
For every $n\geq 2$ and $c\geq 1$ let
\[\vec{v}_c=(c,c) ,\qquad w_c = 1 ,\qquad \vartheta_c = (2c, 1-2c ),\]
and let
$\mathcal{H} (n,c) $ be the irreducible component of $\mathcal{M}(\Lambda_n,\vec{v}_c, 1, \vartheta_c) $
given by the equations
\begin{equation}\label{sliceeqs}
f_1 = f_2 = \cdots = f_{n-1} =0 .
\end{equation}
Then $\operatorname{Hilb}^c (X_n)\simeq \mathcal{H} (n,c) $.
\end{Theorem}
Let $\mbox{pr}\colon \operatorname{Rep}(\Lambda_n, \vec{v}_c, 1)^{\rm ss}_{\vartheta_c} \to \mathcal{M} (\Lambda_n,\vec{v}_c,1, \vartheta_c) $
be the quotient map. The proof {of Theorem~\ref{thm45} basically consists in proving that the counterimage
$\operatorname{pr}^{-1}(\mathcal{H} (n,c) )=:Z_n(c)$ coincides with the total space of the principal fibration $P^n(c)$ we introduced in Section~\ref{background}. As it is quite involved and requires a few intermediate Lemmas and Propositions, we refer the reader to~\cite{bblr} for further details.} Here we only note that the starting point is given by the stability conditions in Definition~\ref{def:stabil}.

\begin{Remark}The set of $(\vec{v}_c,w_c)$-dimensional representations of $\Lambda_n$ which are semistable according to Definition~\ref{def:stabil} does not change if we let the stability parameter vary inside the open cone
\[
 \Gamma_c=\left\{\vartheta=(\vartheta_0,\vartheta_1)\in\mathbb{R}^2\,|\,\vartheta_0>0,\, -\vartheta_0<\vartheta_1<-\frac{c-1}{c}\vartheta_0\right\} .
\]
It can be shown that
for any stability parameter $\bar{\vartheta}$ on the open rays
\begin{gather*}
 R_1 = \big\{(\vartheta_0,\vartheta_1)\in\mathbb{R}^2\,|\,\vartheta_0>0 ,\, \vartheta_0+\vartheta_1=0\big\} ,\\
 R_2 = \big\{(\vartheta_0,\vartheta_1)\in\mathbb{R}^2\,|\,\vartheta_0>0 ,\, (c-1)\vartheta_0+c\vartheta_1=0\big\}
\end{gather*}
there exist representations which are $\bar{\vartheta}$-semistable, but not $\vartheta_c$-semistable.
So, $ \Gamma_c$ is a chamber in the space
$\mathbb{R}^2_{(\vartheta_0,\vartheta_1)}$ of stability parameters and the closed rays $\overline{R_1}$, $\overline{R_2}$ are its walls. Furthermore, inside $\Gamma_c$ semistability and stability are equivalent (cf.~\cite[Lemma~4.7]{bblr}): in particular, points in $\mathcal{M} (\Lambda_n,\vec{v}_c,1, \vartheta_c)$ can be thought of as
$G_{\vec{v}_c}$-orbits of representations in $\operatorname{Rep}(\Lambda_n, \vec{v}_c, 1)$.

A full description of the chamber/wall decomposition of the space $\mathbb{R}^2_{(\vartheta_0,\vartheta_1)}$ will be the object of a future work.
\end{Remark}

We wish to prove that the component $\mathcal{H} (n,c) $ of $\mathcal{M}(\Lambda_n,\vec{v}_c, 1, \vartheta_c) $ introduced in Theorem~\ref{thm45} coincides with the whole moduli space $\mathcal{M}(\Lambda_n,\vec{v}_c, 1, \vartheta_c)$ (this will be Theorem~\ref{mainthm}). Let us introduce the following notation
\begin{equation*}%\label{def_Rnc}
\mathcal{R}(\Lambda_n,c)= \operatorname{Rep}(\Lambda_n,\vec{v}_c,1) ;\qquad \mathcal{R}^{\rm ss}(\Lambda_n,c) =\operatorname{Rep}(\Lambda_n,\vec{v}_c,1)_ {\vartheta_c}^{\rm ss} .
\end{equation*}
Given a representation $(A_1,A_2;C_1, \dots,C_n;e;f_1,\dots,f_{n-1}) \in \mathcal{R}(\Lambda_n,c)$, we form the pencil $A_1 + \lambda A_2$, with $\lambda\in \mathbb{C}$. We recall that a pencil $A_1 + \lambda A_2$ is {\it regular}
if there is a point $[\nu_1,\nu_2]\in \mathbb{P}^1$ such that $\det (\nu_1 A_1 + \nu_2 A_2) \neq 0$.

\begin{figure}[t]\centering
$ \xymatrix@R-2.3em{
&&\mbox{\scriptsize$0$}&&&&\mbox{\scriptsize$1$}\\
&&\bullet\ar@/_3ex/[llddddddd]_{j}\ar@/^/[rrrr]^{a_{1}}\ar@/^4ex/[rrrr]^{a_{2}}&&&&\bullet\ar@/^/[llll]^{b_{1}}
\ar@/^4ex/[llll]^{b_2}\ar@{..}@/^8ex/[llll]\ar@{..}@/^9ex/[llll]
\ar@/^10ex/[llll]^{b_{n-1}} \ar@/^16ex/[llll]^{d_{2}} \ar@/^19ex/[llll]^{d_{3}} \ar@{..}@/^23ex/[llll] \ar@{..}@/^25ex/[llll] \ar@/^28ex/[llll]^{d_{n}} \\ \\ \\ \\ \\ \\ \\
\bullet\ar@/_/[rruuuuuuu]^{i_1}\ar@/_3ex/[rruuuuuuu]_{i_2}\ar@{..}@/_6ex/[rruuuuuuu]\ar@{..}@/_8ex/[rruuuuuuu]\ar@/_10ex/[rruuuuuuu]_{i_{n-1}}&&&&&\\
\mbox{\scriptsize$\infty$}&&&&
}
$
\caption{The quivers $Q'_n$ for $n\ge 3$.}\label{Q'n}
\end{figure}



To prove Theorem \ref{mainthm} it is convenient to introduce ``augmented'' framed quivers defined as follows:
let $Q'_{2}= Q_2$, and, for every $n\geq 3$, let $Q'_{n}$ be the framed quiver in Fig.~\ref{Q'n}. Let $J'_{2}= J_2$ and, for all $n\geq 3$, let $J'_{n}$ be the two sided ideal of $\mathbb{C} Q'_{n}$ generated by the relations
\begin{equation}
\begin{cases}
a_2d_{q+1}-a_1b_q =0, \\
d_{q+1}a_2-b_qa_1-i_qj =0
\end{cases}\qquad\text{for}\quad q=1,\dots,n-1 .
\label{eq-gen-J'}
\end{equation}
We set $\Lambda'_{n}=\mathbb{C} Q'_{n}/J'_{n}$ for all $n\geq2$. Notice that $\Lambda'_{2}=\Lambda_{2}$;
for $n\geq 3$, the algebra $\Lambda_{n}$ can be obtained by taking the quotient of $\Lambda'_{n}$ by a suitable ideal.
Indeed, let $K_{n}$ be the two sided ideal of $\Lambda'_{n}$ generated by the relations
\begin{equation}
\bar{b}_{q}=\bar{d}_q\qquad\text{for}\quad q=2,\dots, n-1 ,\label{eq-gen-K_n}
\end{equation}
where $\bar{x}$ is the class in $\Lambda'_{n}$ of the element $x \in \mathbb{C} Q'_{n}$.
Let $\tilde{p}_{n}\colon\mathbb{C} Q'_{n}\longrightarrow\mathbb{C} Q_{n}$ be the $\mathbb{C}$-algebra morphism determined by the assignments
\begin{equation}\label{eq_tp_n}
\tilde{p}_{n}(a_q) = a_q , \qquad \tilde{p}_{n}(b_q) = c_q ,\qquad \tilde{p}_{n}(d_q) = c_{q} ,\qquad
\tilde{p}_{n}(j) = j , \qquad \tilde{p}_{n}(i_q) = i_q .
\end{equation}
It is straightforward that $\tilde{p}_{n}$ is surjective and that its kernel is the two sided ideal $L_{n}\subset \mathbb{C} Q'_{n}$ generated by the relations
\begin{equation}
b_{q}=d_{q}\qquad\text{for}\quad q=2,\dots,n-1 .
\label{eq-gen-L_n}
\end{equation}
It follows directly from equation~\eqref{eq_tp_n} that $\tilde{p}_n$ maps the set of generators of $J'_n$ (see equation~\eqref{eq-gen-J'}) onto the set of generators of $J_n$ (see equation~\eqref{eq-gen-J}), so that
\begin{equation*}
\tilde{p}_{n}(J'_{n}) =J_{n} .
\end{equation*}
Then it is not hard to check that $\tilde{p}_{n}$ induces a surjective morphism $p_n\colon \Lambda'_n \to \Lambda_n$,
whose kernel, by equations~\eqref{eq-gen-K_n} and \eqref{eq-gen-L_n}, is
\begin{equation*}
\ker p_{n}= L_{n}/(L_{n}\cap J'_{n})=K_{n} .
\end{equation*}
In conclusion, we have proved the following lemma.
\begin{Lemma}\label{lm-Lambda'-->>Lambda}
There is an isomorphism of $\mathbb{C}$-algebras $\Lambda'_{n}/K_{n}\simeq \Lambda_{n}$.
\end{Lemma}

One of the reasons to introduce the augmented quivers $Q'_{n}$ is that their path algebras carry
an action of the group ${\rm SO}(2, \mathbb{C})$ which descends to the quotient algebra $\Lambda'_{n}$. This
action will be instrumental in proving the regularity of the pencil $A_1 + \lambda A_2$.

Elements of ${\rm SO}(2, \mathbb{C})$ will be denoted by $\nu= \left(\begin{smallmatrix} \nu_1 & \nu_2 \\ -\nu_2 & \nu_1\end{smallmatrix}\right)$. Given arrows
\[(a_{1},a_{2};b_{1},\dots,b_{n-1};d_{2},\dots,d_{n};j;i_{1},\dots,i_{n-1})\]
as above and $\nu\in {\rm SO}(2,\mathbb{C})$, we set
\begin{equation*}
\begin{pmatrix}a'_{1} \\ a'_{2} \end{pmatrix} = \nu \begin{pmatrix}a_{1} \\ a_{2} \end{pmatrix},
\qquad
\begin{pmatrix}b'_{q} \\ d'_{q+1} \end{pmatrix} = \nu^{-1}\begin{pmatrix}b_{q} \\ d_{q+1} \end{pmatrix}
\qquad\text{for}\quad q=1,\dots,n-1 .
%\label{eq-a'-->a}
\end{equation*}
The assignment
\begin{gather*}
(a_{1},a_{2};b_{1},\dots,b_{n-1};d_{2},\dots,d_{n};j;i_{1},\dots,i_{n-1}) \\ \qquad{} \longmapsto(a'_{1},a'_{2};b'_{1},\dots,b'_{n-1};d'_{2},\dots,d'_{n};j;i_{1},\dots,i_{n-1}) ,
\end{gather*}
induces an action
\[\widetilde\Phi_n \colon \ {\rm SO}(2,\mathbb{C}) \to \operatorname{Aut}_{\mathbb{C}\hbox{-alg}} (\mathbb{C} Q'_{n}) ,\]
which leaves invariant the generators of the ideal $J'_{n}$, that is,
\begin{equation*}
\widetilde{\Phi}_{n} (\nu) \bigl(J'_{n}\bigr) =J'_{n} .
\end{equation*}
So one has an induced action
\[\Phi_n \colon \ {\rm SO}(2,\mathbb{C}) \to \operatorname{Aut}_{\mathbb{C}\hbox{-alg}} (\Lambda'_{n}) .\]

We wish now to study the space $\operatorname{Rep}(\Lambda'_n, \vec{v}_c, 1) = \mathcal{R}(\Lambda'_n, c)$
of $(c,c)$-dimensional framed representations of $\Lambda'_{n}$ and its open subset
 $\operatorname{Rep}(\Lambda'_n, \vec{v}_c, 1)_{\vartheta_c}^{\rm ss} = \mathcal{R}^{\rm ss}(\Lambda'_n, c)$ of $\vartheta_c$-semistable representations
 (defined analogously to Definition~\ref{def:stabil}).
 For $n=2$ there is nothing new, since
 $\mathcal{R}(\Lambda'_2, c)=\mathcal{R}(\Lambda_2, c)$ and $\mathcal{R}^{\rm ss}(\Lambda'_2, c) = \mathcal{R}^{\rm ss}(\Lambda_2, c)$.
 For $n\geq3$, $\mathcal{R}(\Lambda'_n, c)$ is the affine subvariety of the vector space
\begin{equation*}
\operatorname{Hom}_\mathbb{C}(V_0,V_1)^{\oplus 2}\oplus\operatorname{Hom}_\mathbb{C}(V_1,V_0)^{\oplus 2n-2}\oplus\operatorname{Hom}_\mathbb{C}(V_0,W)\oplus\operatorname{Hom}_\mathbb{C}(W,V_0)^{\oplus n-1}
%\label{eq-point}
\end{equation*}
whose points $(A_1,A_2;B_1,\dots,B_{n-1};D_{2},\dots,D_{n},e;f_1,\dots,f_{n-1})$ satisfy the relations
determined by equations~\eqref{eq-gen-J'}, namely,
\begin{equation}
\begin{cases}
A_2D_{q+1}=A_1B_q,\\
D_{q+1}A_2=B_qA_1+f_{q}e
\end{cases}\qquad\text{for}\quad q=1,\dots,n-1 .
\label{eq-(Q1')}\tag{\rm Q$1'$}
\end{equation}

\begin{Lemma}\label{Q-Lemma}
$\mathcal{R}^{\rm ss}(\Lambda'_{n}, c)$ is the open subset of $\mathcal{R}(\Lambda'_{n}, c)$ determined by the conditions:
\begin{enumerate}\itemsep=0pt
 \item[\rm{(Q$2'$)}]
 for all subrepresentations $S=(S_0,S_1)$ such that $S_0\subseteq\ker e$, one has $\dim S_0\leq\dim S_1$, and, if $\dim S_0=\dim S_1$, then $S=0$;
 \item[\rm{(Q$3'$)}]
 for all subrepresentations $S=(S_0,S_1)$ such that $S_0\supseteq\operatorname{Im} f_i$, for $i=1,\dots,n-1$, one has $\dim S_0\leq\dim S_1$.
\end{enumerate}
\end{Lemma}
\begin{proof}
Given a subrepresentation $(S_{0},S_{1})$, we set $s_{i}=\dim S_{i}$, $i=0,1$. By substituting the definitions of $\vec{v}_c$ and $\vartheta_c$ given in Theorem \ref{thm45} into equations~\eqref{eq:stabnak1} and \eqref{eq:stabnak2} one gets
\begin{gather}
 \text{if $S_0\subseteq \ker e$, then $s_{0}\leq s_1-\dfrac{s_1}{2c}$} ;\label{eq:stabnak1'}\\
\text{if $S_0\supseteq \operatorname{Im} f_i$ for $i=1,\dots,n-1$, then $s_{0}\leq s_{1}+\dfrac{1}{2}-\dfrac{s_{1}}{2c}$} .\label{eq:stabnak2'}
\end{gather}
Whenever $s_1>0$, one has $0<\frac{s_1}{2c}<1$; hence, equation~\eqref{eq:stabnak1'} is equivalent to condition \rm{(Q$2'$)}. On the other hand, as $0\leq\frac{1}{2}-\frac{s_{1}}{2c}<\frac{1}{2}$, equation~\eqref{eq:stabnak2'} is equivalent to condition \rm{(Q$3'$)}.
\end{proof}

\begin{Proposition}\label{prop-regular-in-R'}
For each point of $\mathcal{R}^{\rm ss}(\Lambda'_{n}, c)$ the associated matrix pencil $A_{1}+\lambda A_{2}$ is regular.
\end{Proposition}
\begin{proof}
Let $(A_1,A_2;B_1,\dots,B_{n-1};D_{2},\dots,D_{n},e;f_1,\dots,f_{n-1})$ be a point of $\mathcal{R}^{\rm ss}(\Lambda'_{n}, c)$,
and assume that $A_{1}+\lambda A_{2}$ is singular.
If $c=1$, then $A_{1}+\lambda A_{2}$ is singular if and only if $A_{1}=A_{2}=0$. But this implies the subrepresentation
$(V_0,0)$ does not satisfy condition (Q$3'$). Hence we can assume $c\geq 2$.
The fact that the pencil $A_{1}+\lambda A_{2}$ is singular implies that
there is a nontrivial element
\begin{equation}
v(\lambda)=\sum_{\alpha=0}^{\varepsilon}(-\lambda)^{\alpha}v_{\alpha} \in V_0\otimes_{\mathbb{C}}\mathbb{C}[\lambda]
\label{eq-v(l)}
\end{equation}
such that
\begin{equation}
(A_{1}+\lambda A_{2})v(\lambda)=0\qquad\text{for all $\lambda\in\mathbb{C}$.}
\label{eq-Av(l)}
\end{equation}
By arguing as in the proof of \cite[Lemma~4.11]{bblr}, one can show that the minimal degree
polynomial solution $v(\lambda)$ for the pencil $A_{1}+\lambda A_{2}$ has necessarily degree $\varepsilon >0$.
Let us inductively define the vector spaces $\{U_{i}\}_{i\in \mathbb{N}}$ as follows:
\begin{equation*}
 \begin{cases}
 U_{0}=\langle v_{0},\dots,v_{\varepsilon}\rangle, & \\
 U_{2k+1}=A_{1}(U_{2k})+A_{2}(U_{2k}) &\text{for}\quad k\geq0,\\
 \displaystyle U_{2k}=\sum\limits_{q=1}^{n-1}B_{q}(U_{2k-1})+\sum\limits_{q=2}^{n}D_{q}(U_{2k-1}) &\text{for}\quad k\geq1 .
\end{cases}
\end{equation*}
Note that each $U_{j}$, with $j$ even, is a subspace of $V_0$, while each $U_{j}$, with~$j$ odd, is a subspace of~$V_1$. So, if we introduce the subspaces
\begin{equation*}
S_{0}=\sum_{k=0}^{\infty}U_{2k} \subset V_0 ,\qquad S_{1}=\sum_{k=0}^{\infty}U_{2k+1}\subset V_1 ,
\end{equation*}
it follows that $(S_{0},S_{1})$ is a subrepresentation of $(V_0, V_1)$. We will show that this subrepresentation fails to satisfy either condition~(Q$2'$) or condition~(Q$3'$) of Lemma~\ref{Q-Lemma}, so that one gets a~contradiction.

By substituting equation~\eqref{eq-v(l)} into equation~\eqref{eq-Av(l)} one finds out that
\begin{equation}
 \begin{cases}
A_{1}v_{0} =0,\\
A_{1}v_{\alpha} =A_{2}v_{\alpha-1},\qquad \alpha=1,\dots,\varepsilon,\\
A_{2}v_{\varepsilon} =0,
\end{cases}
\label{eq-A_1v=A_2v}
\end{equation}
so that
\begin{equation}
U_{1}=\langle A_{1}v_{1},\dots,A_{1}v_{\varepsilon} \rangle=A_{1}(U_{0}) .
\label{eq-V_1}
\end{equation}

There are two possible cases, either i) $U_{0}\subseteq\ker e$, or ii) $U_{0}\not\subseteq\ker e$.

i) If we suppose that
$U_{0}\subseteq\ker e$,
equation~\eqref{eq-A_1v=A_2v} and condition \eqref{eq-(Q1')} imply that
\begin{equation*}
U_{2}=\sum_{q=2}^{n}D_{q}A_{1}(U_{0}) .
%\label{eq-V_2-from-V_0}
\end{equation*}
By letting $w_{q,\alpha}=D_{q}A_{1}v_{\alpha}$, $\alpha=1,\dots,\varepsilon$,
for each $q=2,\dots,n$ we obtain an
 element \[(w_{q,1},\dots,w_{q,\varepsilon})\in U_2^{\oplus\varepsilon} \]
 such that $\sum\limits_{\alpha=0}^{\varepsilon-1}(-\lambda)^{\alpha}w_{q,\alpha+1}$ is a polynomial solution for the pencil $A_{1}+\lambda A_{2}$ of degree $\varepsilon - 1$. Since we have supposed $\varepsilon$ to be minimal, one has
 $(w_{q,1},\dots,w_{q,\varepsilon}) = 0$. From that it is easy to deduce that $U_2 =0$ and that $(S_0, S_1) = (U_0,U_1)$.
 So, since $\ker A_1 \cap U_0 \neq 0$ by equation~\eqref{eq-A_1v=A_2v}, then
 equation~\eqref{eq-V_1} entails that $(S_{0},S_{1})$ is a subrepresentation violating condition (Q$2'$).

 ii) Suppose now that $U_{0}$ is not contained in $\ker e$. So, there is at least one $\gamma\in \{0,\dots, \varepsilon\}$ such that
 $e(v_\gamma) \neq 0$. Condition \eqref{eq-(Q1')} implies that
 \begin{equation}
\operatorname{Im} f_{q}= \langle f_{q}e(v_{\gamma})\rangle \subseteq U_{2}\qquad\text{for all}\quad q=1,\dots,n-1 .
\label{eq-imf-in-V}
\end{equation}
To simplify computations, we may assume $\gamma =0$ and $e(v_0)=1$.
 Actually, one checks that the ${\rm SO}(2,\mathbb{C})$ action on $\Lambda'_n$ induces an action on $\mathcal{R}(\Lambda'_{n}, c)$, which commutes with the $G_{\vec{v}_c}$ action defined on the same space, and therefore it restricts to an ${\rm SO}(2,\mathbb{C})$ action on $\mathcal{R}^{\rm ss}(\Lambda'_{n}, c)$. Moreover, this action preserves the regularity of the matrix pencil
 $A_1 + \lambda A_2$. An element $\nu= \left(\begin{smallmatrix} \nu_1 & \nu_2 \\ -\nu_2 & \nu_1\end{smallmatrix}\right)\in {\rm SO}(2,\mathbb{C})$ produces a change of basis
 \[(v_0, \dots, v_{\varepsilon}) \mapsto \nu \cdot (v_0, \dots, v_{\varepsilon}) = (v'_0, \dots, v'_{\varepsilon}) ,\]
 so that
\begin{equation*}
e(v'_{0})=\sum_{\alpha=0}^{\varepsilon}(-\nu_{2})^{\alpha}\nu_{1}^{\varepsilon-\alpha}e(v_{\alpha}) .
\end{equation*}
Since $(e(v_{0}),\dots,e(v_{\varepsilon}))\neq(0,\dots,0)$, there is $\nu\in {\rm SO}(2, \mathbb{C})$ so that $e(v'_{0})\neq0$. Moreover, $e(v'_0)$ can be assumed to be $1$.

Next, by using condition \eqref{eq-(Q1')} and equation~\eqref{eq-A_1v=A_2v}, along with the identity
$A_{2}(\operatorname{Im} f_{q}) =\langle A_{2}f_{q}(1)\rangle$ $ = \langle A_{2}f_{q}(e(v_0))\rangle$, it is not hard to show that
\begin{equation}
A_{2}(\operatorname{Im} f_{q})\subseteq A_{1}(U_{2})\qquad\text{for all}\quad q=1,\dots,n-1 .
\label{eq-A_2(im-f)}
\end{equation}


Now we show that
\begin{equation}
U_{2k+1}\subseteq \sum_{l=1}^{k}A_{1}(U_{2l})
\label{eq-V_(2k+1)-in-imA_{1}}
\end{equation}
for all $k\geq 1$.
Assume $k=1$. By using equations~\eqref{eq-A_1v=A_2v}, \eqref{eq-(Q1')} and condition $e(v_{0})=1$, one gets
\begin{equation*}
f_{q}e(v_{\alpha})=e(v_{\alpha})D_{q+1}A_{1}v_{1}
\end{equation*}
for $q=1,\dots,n-1$. Hence, by using equations~\eqref{eq-A_1v=A_2v} and \eqref{eq-(Q1')} again one shows that
\begin{equation*}
U_{2}= \sum_{q=2}^{n}D_{q}A_{1}(U_{0}) .
%\label{eq-V_2-from-V_0'}
\end{equation*}
 Then $U_{3}$ is spanned by the sets of vectors
\begin{equation*}
\{A_{1}D_{q}A_{1}v_{\alpha}\}_{\begin{subarray}{l} q=2,\dots,n \\ \alpha=1,\dots,\varepsilon \end{subarray}}\subseteq A_{1}(U_{2}) ,\qquad
\{A_{2}D_{q}A_{1}v_{\alpha}\}_{\begin{subarray}{l} q=2,\dots,n \\ \alpha=1,\dots,\varepsilon \end{subarray}}\subseteq A_{2}(U_{2})
%\label{eq-gen-for-V_3}
\end{equation*}
and it follows directly from equations~\eqref{eq-(Q1')} that
$A_{2}D_{q}A_{1}v_{\alpha} \in A_{1}(U_{2})$, for $q=2,\dots,n$. So $U_3\subseteq A_1(U_2)$.

Let us now suppose that equation~\eqref{eq-V_(2k+1)-in-imA_{1}} holds true for $1\leq k \leq m$, with $m\geq 1$. This means that $U_{2m+1}$ is spanned by vectors of the form
$A_{1}w$ with $w\in U_{2l}$, $l=1,\dots,m$.
By noticing that $U_{2m+2}$ is spanned by vectors of the form $B_{p}A_{1}w$ and $D_{q}A_{1}w'$, with $w\in U_{2l}$ and $w'\in U_{2l'}$ for $l,l'=1,\dots,m$, and by using equation~\eqref{eq-(Q1')} and the inductive hypothesis one
finds out that
\begin{equation*}
U_{2m+2}\subseteq \sum_{l=1}^{m} \sum_{q=2}^{n}D_{q}A_{1}(U_{2l})+\sum_{q=1}^{n-1}\operatorname{Im} f_{q} .
%\label{eq-V_(2m+2)-in-imC_qA_{1}+imf_q}
\end{equation*}
From this it follows
\begin{align}
A_1(U_{2m+2}) + A_2(U_{2m+2}) &= U_{2m+3} \nonumber\\
& \subseteq A_{1}(U_{2(m+1)})+\sum_{l=1}^{m} \sum_{q=2}^{n}A_{2}D_{q}A_{1}(U_{2l}) +\sum_{q=1}^{n-1}A_{2}(\operatorname{Im} f_{q}) \nonumber\\
& \subseteq A_{1}(U_{2(m+1)})+A_{1}(U_{2})+\sum_{l=1}^{m} \sum_{q=2}^{n}A_{2}D_{q}A_{1}(U_{2l}) ,\label{last-step}
\end{align}
where in the last step equation~\eqref{eq-A_2(im-f)} has been used.
For $q=2,\dots, n$, equation~\eqref{eq-(Q1')} implies that
\begin{equation*}
A_{2}D_{q}A_{1}(U_{2l}) \subseteq A_{1}(U_{2(l+1)}) .
\end{equation*}
Thus, from equation~\eqref{last-step} we may conclude that
\begin{equation*}
U_{2(m+1)+1}\subseteq\sum_{l=1}^{m+1}A_{1}(U_{2l}) ,
\end{equation*}
so that the inclusion \eqref{eq-V_(2k+1)-in-imA_{1}} is proved.

This and equation~\eqref{eq-V_1} imply that
\begin{equation*}
S_{1}=\sum_{k=0}^{\infty}U_{2k+1}\subseteq A_{1}(U_{0})+\sum_{k=1}^{\infty}\sum_{l=1}^{k}A_{1}(U_{2l})=\sum_{k=0}^{\infty}A_{1}(U_{2k})= A_{1}\left(\sum_{k=0}^{\infty}U_{2k}\right)=A_{1}(S_{0}).
\end{equation*}
But $A_{1}(S_{0})\subseteq S_{1}$, so that
\begin{equation*}
S_{1}=A_{1}(S_{0}) .
\end{equation*}
By equation~\eqref{eq-A_1v=A_2v}, one has $\ker A_{1}\cap S_{0}\neq0$, and therefore $\dim S_{1}<\dim S_{0}$.
Finally, equation~\eqref{eq-imf-in-V} implies that the subrepresentation $(S_{0},S_{1})$ violates condition (Q$3'$).
 \end{proof}

When $n\geq 3$, there is a map $\mathcal{R}(\Lambda_{n}, c) \longrightarrow \mathcal{R}(\Lambda'_{n}, c)$ given by
\begin{equation*}
(A_1,A_2;C_1,\dots,C_n;e;f_1,\dots,f_{n-1}) \mapsto
(A_1,A_2;C_1,\dots,C_{n-1};C_{2},\dots,C_{n},e;f_1,\dots,f_{n-1}) .
\end{equation*}
This map provides a $G_{\vec{v}_{c}}$-equivariant isomorphism of $\mathcal{R}(\Lambda_{n}, c)$ onto the subvariety of $\mathcal{R}(\Lambda'_{n}, c)$ cut by the equations
\begin{equation*}
B_{q}=D_{q}\qquad\text{for}\quad q=2,\dots,n-1
%\label{eq-D_q=C_(q+1)}
\end{equation*}
(cf.\ equations~\eqref{eq-gen-K_n}). Through this isomorphism $\mathcal{R}(\Lambda_{n}, c)$ may be regarded as a closed subvariety of $\mathcal{R}(\Lambda'_{n}, c)$.
\begin{Lemma}\label{lm-R_n(c)-in-R'_n(c)}
When $n\geq 3$, one has that
\begin{equation*}
\mathcal{R}^{\rm ss}(\Lambda_{n}, c) =\mathcal{R}^{\rm ss}(\Lambda'_{n}, c) \cap \mathcal{R}(\Lambda_{n}, c) .
\end{equation*}
\end{Lemma}
\begin{proof}Semistability is a numerical condition which is to be checked on the set of all submodules of a given representation. Hence, it is enough to show that for any left $\Lambda_{n}$-module $M$, an abelian subgroup $N\subset M$ is a left $\Lambda_{n}$-submodule if and only if it is a left $\Lambda'_{n}$-submodule (notice that $M$ has also a natural structure of left $\Lambda'_{n}$-module, induced by restriction of scalars; cf.\ Lemma~\ref{lm-Lambda'-->>Lambda}).
However, precisely because the algebra $\Lambda_n$ is a quotient of $\Lambda'_n$,
the category $\Lambda_{n}$-{\bf mod} is a full subcategory of $\Lambda'_{n}$-{\bf mod}, and this implies in particular that the set of all subobjects of a given $\Lambda_n$-module is the same in the two categories.
\end{proof}

\begin{Theorem}\label{mainthm} The component $\mathcal{H} (n,c) $ of the moduli
space $\mathcal{M}(\Lambda_n,\vec{v}_c, 1, \vartheta_c)$ defined by equations~\eqref{sliceeqs}
coincides with the whole of $\mathcal{M}(\Lambda_n,\vec{v}_c, 1, \vartheta_c)$.
\end{Theorem}
\begin{proof}For each representation $(A_1,A_2;C_1, \dots,C_n;e;f_1,\dots,f_{n-1}) \!\in\! \mathcal{R}^{\rm ss}(\Lambda_n,c)$, equations~\eqref{sliceeqs} hold if and only if the pencil $A_1 + \lambda A_2$ is regular (condition (P2) in Section~\ref{background} and in \cite{bblr}):
\begin{itemize}\itemsep=0pt
 \item in the proof of Proposition~4.9 of~\cite{bblr} it has been shown that condition (P2) holds in $Z_n(c)=\operatorname{pr}^{-1}(\mathcal{H} (n,c))$; i.e., equations~\eqref{sliceeqs} imply the regularity of the pencil;
 \item further on, in the proof of Theorem 4.5 of \cite{bblr} it has been shown that $Z_n(c)$ actually coincides with the open subset of $\mathcal{R}^{\rm ss}(\Lambda_n,c)$ (denoted $\mathcal{R}_n(c)$ in \cite{bblr}) where condition (P2) is satisfied; i.e., the regularity of the pencil implies equations~\eqref{sliceeqs}.
 \end{itemize}
So, the orbit of a $\vartheta_c$-semistable representation \[(A_1,A_2;C_1, \dots,C_n;e;f_1,\dots,f_{n-1}) \in \mathcal{R}^{\rm ss}(\Lambda_n,c)\] lies in $\mathcal{H}(n,c)$ if and only if the pencil $A_1 + \lambda A_2$ is regular. Then the conclusion follows from Proposition \ref{prop-regular-in-R'} and Lemma \ref{lm-R_n(c)-in-R'_n(c)}.
\end{proof}
